\documentclass{standalone}
\usepackage{circuitikz}
\usetikzlibrary{calc}
\definecolor{circuitnotemark}{HTML}{FE00FE}
\begin{document}
\begin{circuitikz}[american, line width=0.8pt]
\ctikzset{bipoles/length=1.2cm}
\ctikzset{grounds/scale=1.36}
\coordinate (e4) at (6,-8);
\coordinate (e6) at (10,-8);
\coordinate (e3) at (4,-8);
\coordinate (b3) at (4,-2);
\coordinate (e7) at (12,-8);
\coordinate (b7) at (12,-2);
\coordinate (h3) at (4,-14);
\coordinate (h7) at (12,-14);
\coordinate (b9) at (16,-2);
\coordinate (h9) at (16,-14);
\coordinate (b11) at (20,-2);
\coordinate (h11) at (20,-14);
\draw (e4) to[esource, l_=$\mathrm{W1}$] (e6); % line 3
\draw (7.82,-8) sin ++(0.09,0.09) cos ++(0.09,-0.09) sin ++(0.09,-0.09) cos ++(0.09,0.09); % line 3
\draw (e3) to[D, l_=$D_{1}$] (b3); % line 4
\draw (e7) to[D, l_=$D_{2}$] (b7); % line 5
\draw (h3) to[D, l_=$D_{3}$] (e3); % line 6
\draw (h7) to[D, l_=$D_{4}$] (e7); % line 7
\draw (b9) to[cC, l_=$C_{1}$, a^=$100\,\mu\mathrm{F}$] (h9); % line 8
\draw (b11) to[R, l_=$R_{L}$, a^=$1\mbox{.}5\,\mathrm{k}\Omega$] (h11); % line 9
\node[ground] at (h9) {}; % line 10
\draw (e3) -- (e4); % line 12
\draw (e6) -- (e7); % line 13
\draw (b3) -- (b7); % line 14
\draw (b7) -- (b9); % line 15
\draw (b9) -- (b11); % line 16
\draw (h3) -- (h7); % line 17
\draw (h7) -- (h9); % line 18
\draw (h9) -- (h11); % line 19
\node[circ] at (e3) {};
\node[circ] at (e7) {};
\node[circ] at (b7) {};
\node[circ] at (h7) {};
\node[circ] at (b9) {};
\node[circ] at (h9) {};
\node[anchor=south west, inner sep=2pt, yshift=4pt, circuitnotemark, font=\LARGE\bfseries] (circuittitle) at (current bounding box.north west) {X};
\path (circuittitle.south west) rectangle ++(10.058,0.853);
\node[anchor=north east, inner sep=2pt, circuitnotemark, font=\large] (circuitstamp) at ([yshift=-0.593cm]current bounding box.south east) {X};
\path (circuitstamp.north east) rectangle ++(-5.08,-0.593);
\end{circuitikz}
\end{document}
